\section{Background}
\label{sec:background}

\subsection{Simulated Annealing}

Simulated Annealing was introduced by \citet{kirkpatrick1983} as a
probabilistic meta-heuristic for combinatorial optimisation, drawing an
analogy with the physical process of slowly cooling a material to reach
a low-energy crystalline state.
The algorithm belongs to the class of \emph{local search} methods: it
maintains a single current solution and perturbs it iteratively.
Unlike greedy local search, SA can accept cost-increasing moves, preventing
premature convergence to poor local optima.

\paragraph{Convergence theory.}
Under a logarithmic cooling schedule $T(t) = c / \log(1+t)$ for sufficiently
large $c$, SA converges to a global optimum with probability
one~\citep{hajek1988}; Kirkpatrick et al.\ \citeyear{kirkpatrick1983}
demonstrated this empirically for combinatorial optimisation.
The logarithmic schedule is impractically slow for large instances.
In practice, geometric schedules---$T(t) = T_0 \cdot r^t$ for some
$r \in (0,1)$---are used.
These converge faster but provide only local-optimum guarantees.
For redistricting bisection, where \metis{} already supplies a good
starting point, local-optimum improvement over a good initial solution
is the operational target.

\paragraph{Geometric cooling.}
The geometric schedule decreases temperature by a fixed multiplicative
factor at each step:
\[
  T(\mathrm{step}) = T_0 \times \left(\frac{T_{\mathrm{final}}}{T_0}\right)^{\mathrm{step}/n_{\mathrm{steps}}}.
\]
This is equivalent to $T_0 \cdot r^{\mathrm{step}}$ with
$r = (T_{\mathrm{final}}/T_0)^{1/n_{\mathrm{steps}}}$.
The schedule spends proportional time at each multiplicative temperature
band, which is natural when the cost landscape has features at many scales.
Linear schedules ($T(t) = T_0 - c \cdot t$) spend proportionally more
time at high temperatures, which may be appropriate for large flat landscapes
but is wasteful for redistricting subgraphs where the \metis{} solution is
already near-optimal.

\subsection{Graph Partitioning and METIS}

The redistricting bisection problem is a variant of graph minimum bisection:
given $G = (V, E, w)$ with $|V| = m$ tracts, find a partition
$(S, V \setminus S)$ with $|S| \approx m/2$ (balanced) that minimises the
weighted edge cut $\EC(S) = \sum_{\{u,v\} \in E: u \in S, v \notin S} w(u,v)$.
The population-balance constraint replaces vertex-count balance, but the
structure is analogous.

\metis{} solves this via multilevel coarsening and KL refinement~\citep{karypis1998}.
Coarsening contracts the graph to a small fraction of its original size by
merging adjacent vertices.
The minimum bisection of the coarsened graph is found by a fast heuristic.
Uncoarsening projects the solution back to the original graph level by level,
applying KL local refinement at each level to improve the solution.
The KL phase swaps boundary vertices between the two parts, accepting swaps
that decrease $\EC$ and stopping when no improving swap exists.
This local-optimum condition is the source of seed sensitivity: a different
coarsening (controlled by the random seed) may reach a different local minimum.

\paragraph{METIS as SA initialisation.}
SA refinement applied \emph{after} \metis{} is computationally sensible
because:
(1) \metis{} provides a feasible, population-balanced, contiguous bisection
    as a warm start;
(2) the \metis{} solution is already near-optimal, so SA needs only fine
    boundary adjustments rather than large structural changes;
(3) \metis{}'s own KL phase exhausts improvements achievable by single-vertex
    swaps with strict acceptance; SA's probabilistic acceptance opens up
    escapes from those local optima.

\subsection{Structure-Layer Algorithms in the B-Series}

The B-series redistricting platform uses a three-layer compositor:
\emph{Structure} (how the bisection tree is built and each node solved),
\emph{WeightsOverride} (what edge weights signal), and
\emph{SeedCompositor} (how many seeds are tried and how the best is selected).
\sa{} is a Structure-layer algorithm.

\paragraph{Existing structure-layer algorithms.}
\geosec{}~\citep{b8paper} constructs the bisection tree optimally with
respect to a ratio objective $\EC/\sqrt{\min(i,k-i)}$, choosing at each level
the split ratio that minimises normalised edge cut.
\areasec{}~\citep{b9paper} adds an area-swing constraint that biases the
geometry of the two pieces.
\ar{}~\citep{b11paper} derives the tree structure from the prime factorisation
of the district count $k$, linking the bisection tree to the Huntington-Hill
apportionment.
All three change \emph{what METIS is optimising} or \emph{how the tree is
structured}, but still delegate the actual bisection solving to a single
\metis{} call per node.

\sa{} differs: it keeps the tree structure (standard or prime-factor, as
configured) and changes the \emph{solver} at each node from a single
\metis{} call to an SA loop seeded by \metis{}.

\paragraph{Orthogonality.}
Because \sa{} is a Structure-layer solver, it composes cleanly with
SeedCompositor strategies.
Running \be{} with $T=50$ seeds, each of which uses \sa{} bisection, would
use $T$ independent SA runs on each node and select the best-EC result.
The empirical results in Section~\ref{sec:empirical} use a single seed with
SA, to isolate the effect of SA refinement from the effect of multi-seed
search.

\subsection{Distinction from ReCom and Short-Burst}

\recom{} is a Markov chain redistricting sampler~\citep{deford2021}.
It proposes plan changes by merging two adjacent districts and re-splitting
them via a spanning tree, accepting all contiguous, balanced proposals.
\recom{}'s goal is to sample from the distribution over valid plans---not
to find the plan with minimum edge cut.
\sa{} is an \emph{optimiser}: it seeks the minimum-$\EC$ bisection at each
tree node and returns a single answer.

Short-burst methods~\citep{cannon2022} run short \recom{} chains from the
current plan, select the best endpoint, and repeat.
This is a different form of optimisation within the Markov chain framework.
\sa{} bisection is structurally simpler: no spanning-tree resampling, no
chain memory, just boundary-tract flips with probabilistic acceptance within
a single bisection node.

The key practical distinction:
\begin{itemize}
\item Use \sa{} (or \be{}) when the goal is a \emph{single best plan}---
  the minimum-edge-cut redistricting for publication or statutory use.
\item Use \recom{} (or ensemble methods) when the goal is
  \emph{distributional coverage}---characterising the space of valid plans
  or auditing an enacted map against the ensemble.
\end{itemize}
